\section{灵敏度分析与一致性检验}
\subsection{已删除对冲扩展的稳定性}
为判断否定对冲是否只是某个随机种子的偶然结果，对已退出正式方案的候选扩展改变情景数、种子与历史池。
表\ref{tab:sens}取自独立灵敏度运行，参考配置为1349.3万元；本轮受控消融的对冲扩展为1349.2万元。
这一0.1万元差异不影响其均高于正式无对冲方案1347.1万元的判断；由于旧灵敏度缺少未舍入逐日记录，不擅自将差异归因为舍入。
\begin{table}[htbp]\centering
\caption{问题三灵敏度运行结果}\label{tab:sens}
\begin{tabular}{lrr}
\toprule
配置 & 费用/万元 & 相对灵敏度参考值\\\midrule
40情景，种子7（参考） & 1349.3 & 0\\
5情景，种子7 & 1350.4 & 0.08\%\\
10情景，种子7/17/27 & 1350.6/1350.4/1351.5 & 0.10/0.08/0.16\%\\
20情景，种子7/17/27 & 1349.2/1350.1/1349.3 & $-0.01$/0.06/0.00\%\\
40情景，种子17/27 & 1350.1/1349.0 & 0.06/$-0.02$\%\\
历史池回看30天 & 1351.3 & 0.15\%\\
历史池回看60天 & 1351.0 & 0.13\%\\
组合权重$\lambda=0.5$ & 1349.1 & $-0.01$\%\\
组合权重$\lambda=0.9$ & 1357.7 & 0.62\%\\
\bottomrule
\end{tabular}
\end{table}
40情景在不同种子下费用变化较小，但各变体并未给出稳定优于1347.1万元正式方案的证据。
短历史池略增费用，说明季节匹配与样本数量需平衡；该表只解释被删除扩展的行为，不再用于给正式策略选情景数。
$\lambda=0.9$明显较差，而0.5与0.7接近，支持保留历史预测成分。

\subsection{跨问题统一参数的代价}
既有问题三参数搜索是在对冲候选结构下完成，不能直接用于宣称无对冲正式方案的独立最优参数。
本轮保持$\lambda=0.7,\kappa=1.02,m=25$不变，只切换是否对冲，以保证消融可归因；其中$\kappa,m$仍沿用问题二开发期选定值，便于跨问比较。
冻结验证期仅用于判断对冲开关，没有反过来重选连续参数，从而避免重复利用验证集。
\subsection{信息与约束的分层检验}
实际执行按式\eqref{eq:discharge}--\eqref{eq:charge}由当前槽构成，结构上不读取未来实现值；残差池检查每个索引严格小于目标日；EWMA只吸收前一日及更早费用。
这些检查验证操作层面的因果性，但全局超参数使用2月至6月开发数据，开发期自身仍是拟合内评价。
另需区别场景LP的完整情景补救和在线因果控制，不能把情景目标值直接当作可实现结算值。

五份结果文件均通过工作表名、日期、行数、计划与紧急量汇总检查。
问题二和三的跨日SOC衔接最大误差为0，问题三计划费及计划量回读差为0；
紧急事件表因小数输出与相邻事件合并产生微小勾稽差，问题二约$2.300\times10^{-4}$ kWh，
问题三约$1.083\times10^{-3}$ kWh，远低于实际调度电量尺度。
